\subsubsection{}

Brukt formel:
\begin{align}
\vec{E} = \frac{\vec{F}}{Q}
\end{align}

som gir
\begin{align}
Q = \frac{\vec{F}}{\vec{E}}
\end{align}

Siden vi kun ser på bevegelse i én dimensjon, kan vi forenkle uttrykket ved å se bort fra vektorformen.

For at partikkelen skal være i ro, må den elektriske kraften balansere tyngdekraften. Det vil si:
\begin{align}
F = -mg
\end{align}

Setter vi dette inn i uttrykket for ladning, får vi:
\begin{align}
Q = \frac{-mg}{E}
\end{align}

Setter inn verdiene:
\begin{align}
Q = \frac{-(0.02\,\text{kg})(-9.81\,\text{m/s}^2)}{100\,\text{N/C}} = 1.962 \cdot 10^{-3}\,\text{C}
\end{align}

Partikkelen må ha motsatt ladning av den øverste platen for at kraften skal virke oppover og balansere tyngdekraften.

\begin{figure}[H]
\centering
\begin{tikzpicture}[>=Stealth, thick]

  % Øverste plate (positiv)
  \fill[gray!30] (-2, 1.5) rectangle (2, 1.8);
  \draw[thick] (-2, 1.5) -- (2, 1.5);
  \foreach \x in {-1.6,-1.0,-0.4,0.2,0.8,1.4} {
    \node at (\x, 1.65) {\small$+$};
  }
  \node[right] at (2.1, 1.65) {$+$};

  % Nederste plate (negativ)
  \fill[gray!30] (-2, -1.8) rectangle (2, -1.5);
  \draw[thick] (-2, -1.5) -- (2, -1.5);
  \foreach \x in {-1.6,-1.0,-0.4,0.2,0.8,1.4} {
    \node at (\x, -1.65) {\small$-$};
  }
  \node[right] at (2.1, -1.65) {$-$};

  % Elektriske feltlinjer
  \foreach \x in {-1.4,-0.7,0.0,0.7,1.4} {
    \draw[->, gray, thin] (\x, 1.45) -- (\x, -1.45);
  }
  \node[right, gray] at (2.3, 0) {$\vec{E}$};

  % Partikkel
  \shade[ball color=blue!60] (0,0) circle (0.18);
  \node[right] at (0.22, 0) {$Q$};

  % Tyngdekraft (nedover)
  \draw[->, red, very thick] (0, -0.18) -- (0, -0.95)
    node[right] {$\vec{G} = m\vec{g}$};

  % Elektrisk kraft (oppover)
  \draw[->, blue!80!black, very thick] (0, 0.18) -- (0, 0.95)
    node[right] {$\vec{F}_Q$};

\end{tikzpicture}
\caption{Partikkelen i ro mellom platene. Den elektriske kraften $\vec{F}_Q$ balanserer tyngdekraften $\vec{G}$.}
\end{figure}

\newpage